%%=============================================================================
%% 摘要
%%=============================================================================

% 中文摘要
\begin{abstract}
运动想象脑机接口（Motor Imagery BCI）技术通过解码运动想象诱发的脑电信号识别用户意图，而时间窗选择是影响分类性能的关键因素。针对现有自适应时间窗优化方法（如深度强化学习）面临的模型复杂度高、训练不稳定、可解释性差等问题，本研究提出一种基于表格型 Q-Learning 的轻量化时间窗优化方法。本研究将时间窗优化问题建模为马尔可夫决策过程，通过维护状态 - 动作价值表迭代估计最优策略，在 BCI Competition IV 2a 数据集上的实验结果表明，所提方法在保持与深度 Q 网络相当分类性能的同时，显著降低了模型复杂度和计算资源消耗。

\keywords{脑机接口；运动想象；时间窗优化；表格型 Q-Learning；轻量化设计}
\end{abstract}

% 英文摘要
\begin{enabstract}
Motor Imagery Brain-Computer Interface (MI-BCI) technology identifies user intentions by decoding EEG signals induced by motor imagery, where time window selection is a critical factor affecting classification performance. To address the limitations of existing adaptive time window optimization methods (such as deep reinforcement learning), including high model complexity, training instability, and poor interpretability, this study proposes a lightweight time window optimization method based on Tabular Q-Learning. The time window optimization problem is formulated as a Markov Decision Process, and the optimal policy is estimated iteratively by maintaining a state-action value table. Experimental results on the BCI Competition IV 2a dataset demonstrate that the proposed method achieves classification performance comparable to Deep Q-Network while significantly reducing model complexity and computational resource consumption.

\enkeywords{Brain-Computer Interface; Motor Imagery; Time Window Optimization; Tabular Q-Learning; Lightweight Design}
\end{enabstract}
